% Generated draft: requires mathematical and source review.
\documentclass[11pt]{article}
\usepackage[margin=25mm]{geometry}
\usepackage{fontspec}
\setmainfont{Latin Modern Roman}
\usepackage{amsmath,amssymb,mathtools,unicode-math}
\setmathfont{Latin Modern Math}
\usepackage{enumitem,needspace,xurl,hyperref}
\setlength{\parindent}{0pt}
\setlength{\parskip}{5pt}
\setlength{\emergencystretch}{3em}
\newcommand{\N}{\mathbb{N}}
\newcommand{\Z}{\mathbb{Z}}
\newcommand{\Q}{\mathbb{Q}}
\newcommand{\R}{\mathbb{R}}
\newcommand{\C}{\mathbb{C}}
\newcommand{\F}{\mathbb{F}}
\newcommand{\A}{\mathbb{A}}
\newcommand{\PP}{\mathbb{P}}
\newcommand{\E}{\mathbb{E}}
\newcommand{\eps}{\varepsilon}
\newcommand{\dd}{\,\mathrm{d}}
\begin{document}
\section*{UnsolvedMath ID: 11000138}
\textbf{Source problem code:} catalogue entry
\subsection{Definitions and mathematical statement}
The referee-contributed question asks what the source's Theorem 5 implies about the end of Riemann's moduli space, especially homology near infinity. For a fixed compact exhaustion \(K_a\) of the intended moduli space \(\mathcal M(F)\), one possible precise end object is the inverse system \(\{H_j(\mathcal M(F)\setminus K_a)\}_a\); the source does not choose this system, its limit or coefficients. Theorem 5 is absent, so no consequence of it is invented.

The following proposed mechanism uses \(F=F_{g,r}^s\), with \(s\geq0\) punctures, and a space \(F^\times\) called its one-point compactification, identifying all punctures. The source says \(F^\times=F\) iff \(s\leq1\); with punctures interpreted as deleted ends, this is not literal equality after one-point compactification, so the marked-point/deleted-end convention must be checked. Let \(C_p(\operatorname{Arc}(F))\) be the source's arc cell complex with boundary \(\partial\). Neither the quotient/cell conventions nor coefficients are provided.

A cell \(\sigma[\alpha]\) comes from an arc family \(\alpha\), and a codimension-one face \(\sigma[\beta]\) removes one arc. The source separates these faces according to whether
\[\operatorname{rank}H_1(F^\times\setminus\!\bigcup\beta)
\quad\text{equals or differs by one from}\quad
\operatorname{rank}H_1(F^\times\setminus\!\bigcup\alpha).\]
The two contributions define \(\partial=\partial_1+\partial_2\), with \(\partial_2\) in the rank-changing case, and are reported to satisfy
\[\partial_1^2=\partial_2^2=0,\qquad
\partial_1\partial_2+\partial_2\partial_1=0.\]
The source's bicomplex has chains on \(p\)-cells in bidegree
\[v=-\operatorname{rank}H_1(F_\alpha),\qquad u=p-v,\]
with \(\partial_1\) of bidegree \((-1,0)\) and \(\partial_2\) of bidegree \((0,-1)\). The cut space \(F_\alpha\) is not defined here or explicitly identified with the complement above. Its stated spectral sequence has initial differential \(\partial_2\) and converges to \(H_*(\operatorname{Arc}(F))\), with convergence/filtration conventions requiring the full construction.

The text explicitly says the following argument is complicated and has not been independently checked, and calls it neither a settled problem nor a theorem: it believes the spectral sequence collapses ``in its \(E^1\)-term to its top horizontal row except in dimension zero.'' It then claims the homology of \(\operatorname{Arc}\) is the \(\partial_1\)-homology of the \(\partial_2\)-kernels in that row. These words could mean concentration at \(E^1\), followed by horizontal homology, rather than literal \(E^1=E^\infty\); the page/degeneration and zero-dimensional exception must be specified. The kernel row must not be replaced by the entire row. Separately it reports [31] that \(\partial_1\)-homology of the top row itself agrees with that of uncompactified \(\mathcal M(F)\).

Finally, the source reports [30,46] that stratifications of arc complexes for bordered surfaces induce a stratification of \(A(F)\) for punctured \(F\), and suggests revisiting calculations [19] and [16,22] while avoiding technical difficulties of the Deligne--Mumford compactification \(\overline{\mathcal M}(F)\). The space \(A(F)\), those calculations and the proposed comparison are not supplied; this is a tentative application, not an established end-homology computation.
\subsection{Short English statement}
Determine what the omitted theorem says about homology near the end of moduli space. The proposed arc-bicomplex mechanism and its top-row spectral-sequence claim are preserved with the source's explicit unchecked-argument warning.
\subsection{Sources}
\begin{itemize}
\item Local UnsolvedMath registry, permanent ID \(11000138\), exact source code \texttt{catalogue entry}.
\item Anonymous referee, contribution in Farb, \emph{Problems on Mapping Class Groups and Related Topics}, Problem 8, PDF page 119: \url{https://www.math.uchicago.edu/~farb/papers/mcgbook.pdf}. Registry access 2026-07-29; pdf-font-structure extraction.
\item Theorem 5; [31] for top-row/moduli homology; [30,46] for stratifications; and [19], [16,22] for proposed calculations are invoked without reproduced statements or full bibliography.
\end{itemize}
\end{document}
